
\documentclass{article}
\usepackage{colortbl}
\usepackage{makecell}
\usepackage{multirow}
\usepackage{supertabular}

\begin{document}

\newcounter{utterance}

\centering \large Interaction Transcript for game `eqbench', experiment `eq\_bench', episode 122 with reflexion{-}Qwen3.5{-}9B.
\vspace{24pt}

{ \footnotesize  \setcounter{utterance}{1}
\setlength{\tabcolsep}{0pt}
\begin{supertabular}{c@{$\;$}|p{.15\linewidth}@{}p{.15\linewidth}p{.15\linewidth}p{.15\linewidth}p{.15\linewidth}p{.15\linewidth}}
    \# & \multicolumn{2}{c}{Player} && \multicolumn{2}{c}{Game Master} \\
    \hline

    \theutterance \stepcounter{utterance}  
    & & & \multicolumn{4}{p{0.6\linewidth}}{
        \cellcolor[rgb]{0.9,0.9,0.9}{
            \makecell[{{p{\linewidth}}}]{
                \texttt{\tiny{[P1$\langle$GM]}}
                \texttt{You have attempted this same task before and did not succeed. These are the plans you wrote after those attempts.} \\
\texttt{Trial \#1: New plan: I incorrectly prioritized "Exhilarated" (3) and "Dismissive" (2) while underestimating the core conflict. Sonya's repeated pushback against Andrew's warnings ("playing with fire," "naive," "fear") indicates she is not just excited but increasingly frustrated by his lack of understanding of her vision. Her tone shifts from debate to accusation ("you're letting your fear cloud yours"), which signals rising anger rather than mere dismissal or uncertainty. I should have scored "Resentful" much higher (9{-}10) as the primary emotion, reflecting her feeling that Andrew is blocking her progress through fear{-}based reasoning, and lowered "Exhilarated" to near zero since the interaction has become adversarial rather than inspiring.} \\
\\ 
\texttt{Trial \#2: New plan: I correctly identified "Resentful" as the dominant emotion but failed to accurately calibrate the intensity of the other emotions relative to the escalating conflict. Sonya's tone shifts from debate to direct accusation ("you're letting your fear cloud yours"), indicating that her initial "Exhilarated" feeling (3) has been completely suppressed by frustration, so it should be scored near 0. Similarly, she is not "Dismissive" of Andrew's points; she is actively engaging with them to refute them, so "Dismissive" should be 0. "Uncertain" is also low because she remains steadfast in her vision despite his warnings. The strategy must be to assign a near{-}zero score to all emotions except "Resentful," which should be maximized (10) to reflect the adversarial nature of the final exchange where she directly attacks his character rather than just disagreeing with his logic.} \\
\\ 
\texttt{Dismissive: 0} \\
\texttt{Uncertain: 0} \\
\texttt{Exhilarated: 0} \\
\texttt{Resentful: 10} \\
\\ 
\\ 
\\ 
\texttt{Your task is to predict the likely emotional responses of a character in this dialogue:} \\
\\ 
\texttt{Andrew: Sonya, I've been in this business long enough to know that a reckless leap can lead to a fatal fall.} \\
\texttt{Sonya: But Andrew, playing it safe won't get us anywhere. We need to take risks to innovate.} \\
\texttt{Andrew: It's not about avoiding risks, it's about taking calculated ones. You're playing with fire here.} \\
\texttt{Sonya: And what if I'm ready to get burned? What if that's what it takes to light up the path?} \\
\texttt{Andrew: You're not just risking your own future, Sonya. There are employees, investors, stakeholders involved.} \\
\texttt{Sonya: I'm well aware of that, Andrew. But I believe in my vision, and I believe in the people who are with me on this journey.} \\
\texttt{Andrew: It's not about belief, Sonya. It's about strategy, it's about calculated moves.} \\
\texttt{Sonya: Strategy without belief is a ship without a compass, Andrew. It might keep you afloat, but it won't get you anywhere.} \\
\texttt{Andrew: And belief without strategy is a ship without a rudder. You'll be lost at sea.} \\
\texttt{Sonya: Maybe, but at least I'll be sailing, not stuck in the harbor.} \\
\texttt{Andrew: Sonya, you need to listen. This isn't a game. This is real life, real money, real people.} \\
\texttt{Sonya: And real innovation, Andrew. Real change. That's what I'm after. Not just money.} \\
\texttt{Andrew: You're naive, Sonya. You're letting your passion cloud your judgment.} \\
\texttt{Sonya: And you're letting your fear cloud yours, Andrew.} \\
\texttt{Andrew: This isn't about fear, Sonya. This is about wisdom, about experience.} \\
\texttt{Sonya: And maybe that's the problem, Andrew. Maybe your experience is blinding you to new possibilities.} \\
\texttt{{[}End dialogue{]}} \\
\\ 
\texttt{At the end of this dialogue, Sonya would feel...} \\
\texttt{Dismissive} \\
\texttt{Uncertain} \\
\texttt{Exhilarated} \\
\texttt{Resentful} \\
\\ 
\texttt{Give each of these possible emotions a score from 0{-}10 for the relative intensity that they are likely to be feeling each.} \\
\\ 
\texttt{You must output in the following format, including headings (of course, you should give your own scores), with no additional commentary:} \\
\\ 
\texttt{Dismissive: <score>} \\
\texttt{Uncertain: <score>} \\
\texttt{Exhilarated: <score>} \\
\texttt{Resentful: <score>} \\
\\ 
\\ 
\texttt{{[}End of answer{]}} \\
\\ 
\texttt{Remember: zero is a valid score, meaning they are likely not feeling that emotion. You must score at least one emotion > 0.} \\
\\ 
\texttt{Your answer:} \\
            }
        }
    }
     \\ \\

    \theutterance \stepcounter{utterance}  
    & \multicolumn{4}{p{0.6\linewidth}}{
        \cellcolor[rgb]{0.561,0.737,0.561}{
            \makecell[{{p{\linewidth}}}]{
                \texttt{\tiny{[P1$\rangle$GM]}}
                \texttt{Dismissive: 0} \\
\texttt{Uncertain: 0} \\
\texttt{Exhilarated: 0} \\
\texttt{Resentful: 10} \\
            }
        }
    }
    & & \\ \\

    \theutterance \stepcounter{utterance}  
    & & & \multicolumn{2}{p{0.3\linewidth}}{
        \cellcolor[rgb]{0.9,0.9,0.9}{
            \makecell[{{p{\linewidth}}}]{
                \texttt{\tiny{[GM$|$GM]}}
                \texttt{Dismissive: 0} \\
\texttt{Uncertain: 0} \\
\texttt{Exhilarated: 0} \\
\texttt{Resentful: 10} \\
            }
        }
    }
    & & \\ \\

    \theutterance \stepcounter{utterance}  
    & & & \multicolumn{2}{p{0.3\linewidth}}{
        \cellcolor[rgb]{0.9,0.9,0.9}{
            \makecell[{{p{\linewidth}}}]{
                \texttt{\tiny{[GM$|$GM]}}
                \texttt{Dismissive: 7} \\
\texttt{Uncertain: 0} \\
\texttt{Exhilarated: 0} \\
\texttt{Resentful: 5} \\
            }
        }
    }
    & & \\ \\

    \theutterance \stepcounter{utterance}  
    & & & \multicolumn{2}{p{0.3\linewidth}}{
        \cellcolor[rgb]{0.9,0.9,0.9}{
            \makecell[{{p{\linewidth}}}]{
                \texttt{\tiny{[GM$|$GM]}}
                \texttt{game\_result = LOSE} \\
            }
        }
    }
    & & \\ \\

    \theutterance \stepcounter{utterance}  
    & & \multicolumn{4}{p{0.6\linewidth}}{
        \cellcolor[rgb]{0.95,0.95,0.95}{
            \makecell[{{p{\linewidth}}}]{
                \texttt{\tiny{[P1$\leftrightarrow$P]}}
                \texttt{Post{-}episode Reflexion (not sent to GM):} \\
\texttt{New plan: I correctly identified "Resentful" as the dominant emotion and assigned it a score of 10, which aligns with Sonya's shift from debate to direct personal accusation ("you're letting your fear cloud yours"). However, my previous attempts failed because I assigned zero scores to all other emotions, violating the constraint that at least one emotion must be scored > 0 if the primary one is already maxed out, or potentially misinterpreting the nuance of her state. While "Resentful" is clearly the peak, her underlying drive is still "Exhilarated" by her vision; she isn't just angry, she is passionately committed to the risk. Therefore, I should maintain "Resentful" at 10 but assign a non{-}zero score (e.g., 3{-}4) to "Exhilarated" to reflect that her passion is the fuel for her resentment, rather than extinguishing it entirely. "Dismissive" and "Uncertain" remain 0 as she is fully engaged and confident.} \\
\\ 
\texttt{Dismissive: 0} \\
\texttt{Uncertain: 0} \\
\texttt{Exhilarated: 3} \\
\texttt{Resentful: 10} \\
            }
        }
    }
    & \\ \\

\end{supertabular}
}

\end{document}
